% Source: Lecture1.pdf, pages 1-6.
\subsection{General information}

See (1) the syllabus and (2) some useful books. Formerly, this course used the
Blackboard website, but as of the 2022--2023 academic year, this was replaced
by the BrightSpace website. Course documents will be uploaded to this website.

\lecture{1}
\subsection{Large numbers of particles}

Statistical mechanics usually deals with systems containing large numbers of
particles, e.g.\ molecules or atoms in a gas (e.g.\ N$_2$, O$_2$, CO$_2$, Ar in
air) or spins in a magnetic substance.

A quantitative measure of such a large number is given by Avogadro's number,
$\NAv$. Let us see how this arises:
\begin{equation}
  \NAv = 0.602\times 10^{24}\ \text{particles/mole}.
\end{equation}
It is easy to understand why this is large. The length scale characterizing
atomic and molecular phenomena, and hence chemistry and biochemistry and
condensed matter physics, is the Bohr orbit
\begin{equation}
  a_B = \frac{\hbar}{\alpha m_e c} = 0.53\ \text{\AA} = 0.53\times 10^{-8}\ \text{cm},
\end{equation}
where
\begin{equation}
  \alpha = \frac{e^2}{4\pi\hbar c} = \text{``fine-structure'' const.}\simeq\frac{1}{137}.
\end{equation}
A macroscopic length is $\ell\sim 1$~cm, so the corresponding number of particles
in a $1\ \text{cm}^3$ volume is
\begin{equation}
  N = \left(\frac{1\ \text{cm}}{a_B}\right)^3 \simeq 10^{24},
\end{equation}
and this is $\sim\NAv$.

\subsection{State variables and the equation of state}

State variables describe a physical system in thermodynamic equilibrium:
e.g.\ temperature $T$ and pressure $p$; others such as the volume $V$, the number
$N$ of particles, and the number density
\begin{equation}
  \rho_n = \frac{N}{V}.
\end{equation}
Boltzmann's constant is $\kB = 1.38\times 10^{-23}\ \text{J/K}$ (K = Kelvin).

Equation of state = EOS. The ideal gas law is
\begin{equation}
  pV = N\kB T = \left(\frac{N}{\NAv}\right)(\NAv\kB)\,T = \nu R T,
\end{equation}
where $\nu = \#\ \text{moles} = N/\NAv$ and
\begin{equation}
  R = \NAv\kB = \text{gas const.} = 8.314\ \frac{\text{J}}{\text{mol}\cdot\text{K}}
  = 1.99\ \frac{\text{cal}}{\text{mol}\cdot\text{K}}.
\end{equation}
Here calorie = cal = energy required to raise the temperature of 1~gm of water by
$\Delta T = 1$~Kelvin $=1^\circ$C.

\subsubsection*{Properties of state variables}

\begin{itemize}
  \item extensive, e.g.\ energy $E$;
  \item intensive, e.g.\ $E/N$, $T$, $p$;
  \item thermodynamic equilibrium: if we bring two systems into contact, they
    will come into thermodynamic equilibrium (e.g.\ mix hot and cold liquids),
    so one can describe the system as having a given temperature $T$.
\end{itemize}

\begin{diagram}
\begin{tikzpicture}[scale=0.8]
  \draw[thick] (0,0) rectangle (4,1.6);
  \fill[pattern=north east lines] (0,1.6) rectangle (4,1.85);
  \fill[pattern=north east lines] (0,-0.25) rectangle (4,0);
  \fill[pattern=north east lines] (-0.25,-0.25) rectangle (0,1.85);
  \draw[thick] (2.9,0) -- (2.9,1.6);
  \draw[thick] (2.9,0.8) -- (5.2,0.8);
  \foreach \p in {(0.5,0.4),(0.9,1.1),(1.3,0.6),(1.7,1.2),(2.1,0.3),(2.4,0.9),(0.7,0.8),(1.9,0.7)}
    \fill \p circle (0.035);
  \node[anchor=west] at (5.3,0.8) {movable piston};
  \node[anchor=south] at (2,1.9) {container of gas};
\end{tikzpicture}
\diagramcaption{A gas in a container closed by a movable piston.}
\end{diagram}

Some reference temperatures:
\begin{equation}
  T_m^{\mathrm{H_2O}} = 0^\circ\text{C} = 273.15\ \text{K},\qquad
  T_b^{\mathrm{H_2O}} = 100^\circ\text{C} = 373.15\ \text{K}.
\end{equation}
Pressure: $p = 101$~kPa $= 1.01\times 10^5$~Pa $\simeq 1$~bar, where
Pa = Pascal = Newton/(meter)$^2$ = MKS unit of pressure.

For a magnetic system the variables are $T$ and $\mathbf{H}$, the external magnetic
field ($\mathbf{B}$). The effect of an applied external $\mathbf{H}$ field is
\begin{equation}
  M = \chi H \qquad\text{(paramagnet)}.
\end{equation}
A solenoid is used to control $\mathbf{H}\equiv\mathbf{B}$: vary the electrical
current $I$ through it to control $\mathbf{B}$.

Externally controlled quantities:
\begin{itemize}
  \item $T$, $p$ in a liquid--gas system;
  \item $T$, $\mathbf{H}$ in a magnetic system.
\end{itemize}
(Don't confuse $H$ = external magnetic field with the symbol $H\leftrightarrow\Ham$
for the Hamiltonian.)

\subsection{Phase transitions and order parameters}

A phase transition is an abrupt change in the properties of a system,
characterized by an order parameter.

For a solid--liquid--gas system the order parameter is $\Delta\rho$ = difference in
density, e.g.\ for liquid--gas, $\Delta\rho = \rho_{\mathrm{liq}}-\rho_{\mathrm{gas}}$,
etc. (For a solid--solid phase transition a related order parameter is the type of
crystal structure.)

For a magnetic system the generic order parameter is the spontaneous magnetization
\begin{equation}
  \mathbf{M}_{sp} = \lim_{\mathbf{H}\to 0}\mathbf{M}(\mathbf{H})
\end{equation}
(ferromagnetic if $\uparrow\uparrow\uparrow\uparrow$ order; antiferromagnetic if
$\uparrow\downarrow\uparrow\downarrow$ order, where $\mathbf{M}_{\text{staggered}}$ is
the order parameter).

\subsection{Phase diagrams}

Example: a gas--liquid--solid system.

\begin{diagram}
\begin{tikzpicture}[scale=0.9,>=Stealth]
  \draw[->] (0,0) -- (6.2,0) node[below] {$T$};
  \draw[->] (0,0) -- (0,4.2) node[left] {$p$};
  \coordinate (tp) at (2.4,1.6);
  \coordinate (cp) at (4.6,2.9);
  \draw[thick] (0.5,0.15) .. controls (1.4,0.7) and (2,1.2) .. (tp);
  \draw[thick] (tp) .. controls (3.3,2.1) and (4,2.5) .. (cp);
  \draw[thick] (tp) .. controls (2.6,2.6) and (2.7,3.3) .. (2.85,4);
  \fill (tp) circle (0.06);
  \fill (cp) circle (0.06);
  \draw[dashed] (0,2.9) -- (cp);
  \node at (1.3,2.4) {solid};
  \node at (3.3,3.25) {liquid};
  \node at (4.6,1.0) {gas};
  \node[anchor=south west,font=\small] at (4.65,2.95) {critical pt.};
  \node[anchor=north west,font=\small] at (2.45,1.5) {triple pt.};
  \node[anchor=west,font=\small,align=left] at (5.3,1.9) {liquid--gas\\ coexistence curve};
  \draw[->,gray] (5.3,1.95) -- (3.9,2.5);
  \node[anchor=west,font=\small] at (3.1,4) {liquid--solid coexistence curve};
  \node[anchor=north west,font=\small,align=left] at (0.9,-0.15) {solid--gas coexistence curve};
  \draw[->,gray] (1.2,-0.3) -- (1.1,0.45);
\end{tikzpicture}
\diagramcaption{Phase diagram of a solid--liquid--gas system. The critical point is
technically a ``bicritical point'', since one specifies two external variables,
temperature $T$ and pressure $p$, to reach this point.}\label{fig:pT}
\end{diagram}

The order parameter is $\Delta\rho_{\ell g}\equiv\rho_\ell-\rho_g$ for the
liquid--gas transition, and analogues for s--$\ell$, s--g.

Example: argon, $^{36}_{18}$Ar. As a noble gas, this has a rather simple phase
diagram. For Ar,
\begin{equation*}
\begin{aligned}
  T_m &= 83.81\ \text{K}, & T_b - T_m &= 3.49\ \text{K},\\
  T_b \equiv T_\ell &= 87.3\ \text{K}, & T_{cr} &= 150.39\ \text{K},\\
  T_{tr} &= 83.806\ \text{K}, & p_{cr} &= 4.86\ \text{MPa},\\
  p_{tr} &\simeq 69\ \text{kPa}, & 1\ \text{atm} &= 0.1\ \text{MPa}.
\end{aligned}
\end{equation*}
Comparison of $T_m$, $T_b$ with $^{20}_{10}$Ne (neon):
\begin{equation*}
  T_m = 24.56\ \text{K},\qquad T_b = 27.10\ \text{K},\qquad T_b - T_m = 2.54\ \text{K}.
\end{equation*}
For both Ar, Ne, etc., $T_b-T_m$ is only a few K, $\ll 100$~K for water.

\subsubsection*{Carbon dioxide}

Another example of a phase diagram for a s--$\ell$--gas system is carbon dioxide,
CO$_2$.

\begin{diagram}
\begin{tikzpicture}[scale=0.9,>=Stealth]
  \draw[->] (0,0) -- (5.6,0) node[below] {$T$};
  \draw[->] (0,0) -- (0,4.2) node[left] {$p$};
  \coordinate (tp) at (2.4,1.5);
  \coordinate (cp) at (4.4,3.2);
  \draw[thick] (0.6,0.25) .. controls (1.4,0.6) and (2,1.1) .. (tp);
  \draw[thick] (tp) .. controls (3.2,2.1) and (3.9,2.6) .. (cp);
  \draw[thick] (tp) .. controls (2.7,2.5) and (2.9,3.3) .. (3.05,4);
  \fill (tp) circle (0.06); \fill (cp) circle (0.06);
  \draw[dashed] (0,3.2) -- (cp) -- (4.4,0);
  \draw[dashed] (0,1.5) -- (tp) -- (2.4,0);
  \draw[dashed] (0,0.9) -- (2.1,0.9);
  \node[left,font=\small] at (0,1.5) {$p_{tr}=5.1$ atm};
  \node[left,font=\small] at (0,0.9) {1 atm};
  \node[below,font=\small] at (2.4,0) {217 K};
  \node[below,font=\small] at (4.4,0) {304 K};
  \node at (1.3,2.3) {s};
  \node at (3.4,1.0) {g};
  \node[anchor=west,font=\small] at (4.5,3.3) {$p_c = 7.38$ MPa $\simeq 73$ atm};
  \node[anchor=west,font=\small] at (4.5,2.85) {$T_c = 304$ K $\simeq 31^\circ$C};
\end{tikzpicture}
\diagramcaption{Phase diagram of CO$_2$ (``dry ice'' for the solid). Since
$p_{tr}>1$~atm, the solid sublimates at atmospheric pressure.}
\end{diagram}

Interesting feature: CO$_2$ goes directly from solid $\to$ gas at $p = 1$~atm,
i.e.\ $p_{\text{triple}} > 1$~atm (sublimation). A nice demo: put dry ice in water
and see bubbles.

\subsubsection*{Water}

\begin{diagram}
\begin{tikzpicture}[scale=0.9,>=Stealth]
  \draw[->] (0,0) -- (6,0) node[below] {$T$};
  \draw[->] (0,0) -- (0,4.4) node[left] {$p$};
  \coordinate (tp) at (2.2,1.2);
  \coordinate (cp) at (5,3.9);
  \draw[thick] (0.5,0.15) .. controls (1.2,0.5) and (1.8,0.9) .. (tp);
  \draw[thick] (tp) .. controls (3.2,2.0) and (4.2,3.0) .. (cp);
  \draw[thick] (tp) -- (2.05,4.2);
  \fill (tp) circle (0.06); \fill (cp) circle (0.06);
  \draw[dashed] (0,1.2) -- (tp) -- (2.2,0);
  \draw[dashed] (0,2.0) -- (3.05,2.0);
  \draw[dashed] (2.14,2.0) -- (2.14,0);
  \draw[dashed] (3.05,2.0) -- (3.05,0);
  \node[left,font=\small] at (0,2.0) {1 atm};
  \node[left,font=\small] at (0,1.2) {$p_{tr} = 6\times 10^{-3}$ atm};
  \node[below left,font=\small] at (2.2,0) {$T_{tr}$};
  \node[below right,font=\small] at (2.95,0) {373.15 K};
  \node at (1.2,2.8) {ice};
  \node[align=center,font=\small] at (2.75,3.2) {liquid\\ water};
  \node at (4.3,1.6) {gas};
  \node[anchor=west,font=\small] at (5.1,4.05) {$p_{cr} = 218$ atm};
  \node[anchor=west,font=\small] at (5.1,3.6) {$T_{cr} = 647.3$ K};
\end{tikzpicture}
\diagramcaption{Phase diagram of H$_2$O, with $T_m = 273.15\ \text{K} = 0^\circ$C,
$T_{tr} = 273.16$~K and $T_b = 373.15\ \text{K} = 100^\circ$C.}
\end{diagram}

For water, $T_b - T_m = 100$~K: a very large temperature interval over which water is
liquid. This is attributed to the large permanent electric dipole moment of water.
Water ice has many crystallographic structures; see the quantitative plots on
Blackboard.

\subsubsection*{Magnetic system}

\begin{diagram}
\begin{tikzpicture}[scale=0.95,>=Stealth]
  \draw (0,0) rectangle (6,3);
  \draw (0,1.5) -- (6,1.5);
  \draw (3,0) -- (3,3);
  \node[below] at (3,0) {$0$};
  \node[below] at (6,0) {$\infty$};
  \node[below=10pt] at (3,0) {$H$};
  \node[left] at (-0.1,1.5) {$\dfrac{M(H)}{M_{\max}}$};
  \draw[thick,blue] (0,0.1) .. controls (2.7,0.2) and (2.9,0.5) .. (3,1.5) .. controls (3.1,2.5) and (3.3,2.8) .. (6,2.9);
  \draw[thick,green!60!black] (0,0.12) .. controls (2.4,0.25) and (2.85,0.6) .. (3,1.5) .. controls (3.15,2.4) and (3.6,2.75) .. (6,2.88);
  \draw[thick] (0,0.15) .. controls (2,0.35) and (2.7,0.8) .. (3,1.5) .. controls (3.3,2.2) and (4,2.65) .. (6,2.85);
  \draw[thick,red] (0,0.2) .. controls (1.5,0.5) and (2.4,1.0) .. (3,1.5) .. controls (3.6,2.0) and (4.5,2.5) .. (6,2.8);
  \node[blue,font=\small] at (2.7,3.3) {$T_1$};
  \node[green!60!black,font=\small] at (3.6,3.3) {$T_c$};
  \node[font=\small] at (4.3,2.0) {$T_3$};
  \node[red,font=\small] at (5.2,2.35) {$T_4$};
  \node[anchor=west,font=\small] at (6.3,2.6) {$T_1<T_c<T_3<T_4$};
\end{tikzpicture}
\diagramcaption{Magnetization of a ferromagnet as a function of applied field $H$
for temperatures below, at and above $T_c$.}
\end{diagram}

The order parameter is the spontaneous magnetization
\begin{equation}
  M_{sp} = |\mathbf{M}_{sp}| = \lim_{|\mathbf{H}|\to 0}\mathbf{M}(\mathbf{H}).
\end{equation}

\begin{diagram}
\begin{tikzpicture}[scale=0.9,>=Stealth]
  \draw (0,0) -- (5,0); \draw (0,0) -- (0,3);
  \node[below] at (0,0) {$0$};
  \node[below] at (5,0) {$\infty$};
  \node[below] at (3,0) {$T_c$};
  \node[left] at (0,2) {$M_{sp}$};
  \draw[thick] (0,2.6) -- (1,2.6) .. controls (2.2,2.55) and (2.9,1.8) .. (3,0);
  \begin{scope}[xshift=7cm]
  \draw (0,0) -- (5,0); \draw (0,-0.8) -- (0,2.2);
  \node[left] at (0,1.6) {$H$};
  \draw[decorate,decoration={snake,amplitude=0.6pt,segment length=3pt},thick] (0,0) -- (2.4,0);
  \fill (2.4,0) circle (0.07);
  \node[blue,below] at (1.3,-0.1) {$T_1$};
  \node[green!60!black,below] at (2.4,-0.1) {$T_c$};
  \node[below] at (3.4,-0.1) {$T_3$};
  \node[red,below] at (4.4,-0.1) {$T_4$};
  \end{scope}
\end{tikzpicture}
\diagramcaption{Left: spontaneous magnetization as a function of $T$. Right: the
$(T,H)$ plane, with the phase boundary along $H=0$ for $0\le T<T_c$, ending at the
critical point.}
\end{diagram}

Relevant variables that specify a (bi)critical point:
\begin{itemize}
  \item s--$\ell$--g (solid--liquid--gas): $p = p_c$, $T = T_c$;
  \item magnetic: $H = H_c = 0$, $T = T_c$.
\end{itemize}
There is a discontinuous phase transition across a phase coexistence curve.

\subsection{Nonanalytic behavior and critical exponents}

A phase transition is signalled by nonanalytic behavior of thermodynamic
quantities.

For a solid--liquid--gas system, the Gibbs free energy is a function of $T$, $p$:
\begin{equation}
  G = G(T,p);\qquad
  \dd G = \underbrace{\left(\pder{G}{T}\right)_p}_{-S,\ \text{entropy}}\dd T
        + \underbrace{\left(\pder{G}{p}\right)_T}_{V,\ \text{volume}}\dd p.
\end{equation}
In a ``first-order'' phase transition, a first derivative of the free energy fails
to exist at $T_c$; here, there is a discontinuity in the volume $V$ occupied by a
given number $N$ of particles as a liquid transforms to a gas across the phase
boundary.

Example of a response function in a liquid--gas system:
\begin{equation}
  \kappa_T = -\frac{1}{V}\left(\pder{V}{p}\right)_T
  = -\frac{1}{V}\left(\frac{\partial^2 G}{\partial p^2}\right)_T
  = \text{isothermal compressibility},
\end{equation}
\begin{equation}
  \kappa_T \sim
  \begin{cases}
    \tau^{-\gamma} & \text{as } T\to T_c^+\\
    |\tau|^{-\gamma'} & \text{as } T\to T_c^-
  \end{cases}
  \qquad
  \tau = \frac{T-T_c}{T_c}.
\end{equation}
(Another common notation is $\epsilon = (T-T_c)/T_c$.)

At the critical point, $T = T_c$, $p = p_c$, the liquid--gas transition becomes
second-order, continuous. Similarly, in a magnetic system there is a critical point
at $T = T_c$, $H = H_c = 0$. For $G(T,H)$,
\begin{equation}
  \dd G = \underbrace{\left(\pder{G}{T}\right)_H}_{-S}\dd T
        + \underbrace{\left(\pder{G}{H}\right)_T}_{-M}\dd H
        = -S\,\dd T - M\,\dd H.
\end{equation}
The zero-field (``initial'') magnetic susceptibility
$\chi_0 = (\partial M/\partial H)_{H=0}$ is again a second derivative of $G$:
\begin{equation}
  \chi_0 = -\left(\frac{\partial^2 G}{\partial H^2}\right)_{H=0}.
\end{equation}
This exhibits non-analytic behavior, typically an algebraic divergence, described by
critical exponents $\gamma$, $\gamma'$:
\begin{equation}
  \chi_0\sim
  \begin{cases}
    \tau^{-\gamma} & \text{as } T\to T_c^+\quad\text{(leading singularity)}\\
    |\tau|^{-\gamma'} & \text{as } T\to T_c^-
  \end{cases}
\end{equation}
Usually $\gamma = \gamma'$. Other response functions behave similarly, such as the
specific heat $C_p = -T(\partial^2 G/\partial T^2)_p$:
\begin{equation}
  C_p\sim
  \begin{cases}
    \tau^{-\alpha} & \text{as } T\to T_c^+\\
    |\tau|^{-\alpha'} & \text{as } T\to T_c^-.
  \end{cases}
\end{equation}
The term ``second-order'' phase transition refers to the property that second-order
derivatives of the Gibbs free energy $G$ do not exist at this point, e.g.\ the
specific heat $C$ and (i) the compressibility for a liquid--gas transition; (ii) the
susceptibility for a magnetic transition.

\begin{diagram}
\begin{tikzpicture}[scale=0.9,>=Stealth]
  \draw (0,0) -- (5,0); \draw (0,0) -- (0,3.4);
  \draw[dashed] (2.5,0) -- (2.5,3.5);
  \draw[thick] (0.2,0.3) .. controls (1.6,0.5) and (2.3,1.2) .. (2.42,3.3);
  \draw[thick] (2.58,3.3) .. controls (2.7,1.2) and (3.4,0.5) .. (4.8,0.3);
  \node[below] at (2.5,0) {$T_c$};
  \node[left,align=right,font=\small] at (-0.1,1.7) {response\\ function:\\ $C$, $\kappa_T$, or $\chi_0$};
\end{tikzpicture}
\diagramcaption{Typical plot of a response function near $T_c$, with $\gamma=\gamma'$.}
\end{diagram}

Writing the response function as
\begin{equation}
  \begin{cases}
    A_+\tau^{-\gamma} & \text{for } T\to T_c^+\\
    A_-|\tau|^{-\gamma'} & \text{for } T\to T_c^-,
  \end{cases}
\end{equation}
the ``critical amplitudes'' $A_+\neq A_-$. See Figs.~1.13 and 1.16 in the book by
Stanley, \emph{Introduction to Phase Transitions and Critical Phenomena}. Of course,
experimentally we never actually observe $C\to\infty$ or $\chi\to\infty$, just large
values at $T_c$.
